\documentclass{article}
\usepackage{graphicx} % Required for inserting images
\usepackage{amsmath}
\usepackage{amssymb}
\usepackage{algorithm}
\usepackage{algpseudocode}
\usepackage{hyperref}
\usepackage{booktabs}
\usepackage{float}
\usepackage{url}       % For \url{}

\title{FinalProjectReport\_NLP}
\author{Trent Hair, Jion Iwasaki, Tommy Le}
\date{May 2025}

\begin{document}

\maketitle

\begin{abstract}
This paper presents an enhanced version of the Parallel Temporal Neural Coding Network (PTNCN) architecture, incorporating fast weights and gate mechanisms for improved sequential data processing. Through extensive experimentation on character-level language modeling tasks, we demonstrate that while the original PTNCN achieves higher accuracy (32.77\%), the enhanced PTNCN-FW model offers a more robust solution with competitive accuracy (28.38\%) and significantly better training stability. The integration of fast weights and gate mechanisms provides better control over information flow and parameter growth, making the model more suitable for real-world applications. Our comparative analysis with LSTM models shows that the enhanced architecture outperforms traditional approaches in both accuracy and training stability, while maintaining efficient resource utilization. The findings highlight the importance of balancing performance and stability in neural network design, offering valuable insights for future research in sequential learning architectures.

\end{abstract}

\section{Introduction}
\hspace*{.5cm}Many Language Learning Models rely heavily on recurring neural networks and backpropagation to recognize patterns as a sequence. Backpropagation (BPTT) uses the chain rule of the derivatives to approximate how all the weights affect the overall loss. Based on this value, we can adjust all the weight to reduce the loss and achieve gain by moving in the opposite direction of the gradient. 

Although BPTT has been the standard for training neural networks, BPTT does not permit the use of non-differentiable activation functions. All operations, including the activation function, must be differentiable (smooth, continuous) to compute the gradient. RNN trained with BPTT struggle to handle sequential data streams, such as real-time speech recognition or stock market prediction. For sequences with longer input range, they require creating larger computational graphs, which could lead to vanishing or exploding gradients. Moreover, there is a credit assignment issue, in which BPTT calculates global loss by accumulating contributions for all hidden layers for the weights. If we were to troubleshoot, we wouldn't be able to identify faulty layers. Finally, BPTT is sequential by nature, making it computationally inefficient as it can’t take advantage of massively parallel systems.

Parallel Temporal Neural Coding Networks (P-TNCN) were based on Bayesian Brain theory or Predictive Coding, which models the brain as a system that generates predictions about incoming data and continuously adjusts these predictions based on the observed outcomes. The P-TNCN architecture aims to overcome several limitations of traditional RNNs trained with BPTT. Specifically, it addresses issues related to non-differentiable activation functions, computational inefficiency, and difficulty in handling long-term dependencies.

Unlike BPTT, which calculates a global error signal and back-propagates it through all network layers and time stamps, P-TNCN adopts Local Representation Alignment (LRA). This approach enables each layer to independently align its local representations, significantly reducing the reliance on sequential computations. The model achieves this by utilizing a hierarchical structure of recurrent layers, each of which generates its predictions and corrects them using local error signals. This local correction mechanism allows P-TNCN to operate efficiently even with non-differentiable activation functions, making it suitable for a broader range of sequence modeling tasks.

Fast Weights are an enhancement to the standard P-TNCN architecture, inspired by the concept of synaptic plasticity in biological neural networks. Traditional neural networks rely on slow, gradually adjusted weights that store long-term knowledge. In contrast, Fast Weights provide a rapidly adjustable memory mechanism that can store and recall information from recent inputs, effectively acting as a short-term memory. This mechanism is particularly beneficial for sequence processing tasks where recent information is crucial for making accurate predictions.

In our work, we integrated Fast Weights into the P-TNCN architecture to enhance the ability to capture short-term dependencies in sequential data. This integration allows the model to quickly adapt to changing patterns in the input sequence without affecting the stability of the long-term weights. By carefully tracking the evolution of both Fast Weights and regular weights during training, we were able to gain insights into the model’s learning dynamic and demonstrate the advantages of combining Fast Weights with the P-TNCN framework.


\section{Dataset}
\subsection{- Dataset Description}
We use the Penn TreeBank (PTB) dataset, specifically the character-level version, for our experiments. The dataset is split into three parts:
\begin{itemize}
    \item Training set (trainX.txt): containing the main training data
    \item Validation set (validX.txt): used for model validation
    \item Test set (testX.txt): used for final evaluation
    \item Subset (subX.txt): a smaller subset used for quick experiments
\end{itemize}

\subsection{- Dataset Preprocessing}
The dataset consists of character-level sequences with a vocabulary of 50 characters, including:
\begin{itemize}
    \item Lowercase letters (a-z)
    \item Numbers (0-9)
    \item Special characters (\#, \$, ', \&, *, -, /, ., <, >, \_, \textbackslash)
\end{itemize}

The preprocessing steps include:
\begin{enumerate}
    \item Character-level tokenization of the input text
    \item Conversion of characters to one-hot encoded vectors
    \item Sequence padding to handle variable-length inputs
    \item Creation of input-target pairs for training
\end{enumerate}

The data is processed in batches to facilitate efficient training, with each batch containing multiple sequences. The model is trained to predict the next character in the sequence, making it a character-level language modeling task.

\section{Model Architecture: PTNCN with Gate and Fast Weight Mechanisms}
% --- FOrward enndpoint
\subsection{Forward Pass}
\subsubsection{Forward pass with gates}
\hspace*{.5cm} Similar to how the reset gate is applied in GRUs, our gate mechanism was designed based on the forward pass equation of the original P-TNCN model. Its size was determined to reflect both the influence from lower layers and the previous state of the current layer.

Top-down feedback is used to guide or correct lower layers, but "Forward Passing" the layers primarily relies on its own past activity (self-memory) and bottom-up input to make initial predictions.
It does not require top-down feedback to compute its output, only to refine it later based on what the higher layer (Layer 2) expects.
In early computation, Layer 1 should focus on aligning with the bottom-up signal rather than trying to anticipate the top-down correction.

Forward Passing also incorporates Fast weights. Fast weights matrix represent the connection between the previous state and the current state. By multiplying the fast weight matrix with the previous activations, the model reinforces patterns it has recently experienced and revives important signals into the current state. 


\begin{itemize}
    \item $G^l_1, G^l_2$, : the two gate weight needed to update 
    \item $V_l$, : weights that incorporate the bottom up signal. The purpose is so the output matches with the lower layers. 
    \item $U_l$, :  weights that incorporate the previous time step of same layer 
    \item $y_l^{t-1}$:
    \item $A^t_l$: \textbf{Fast weights to reinforce connection between the pass and current activation}
    \item $\phi_z^l$: activation function for the forward pass
\end{itemize}
\begin{align}
g_l^t &= \sigma(G_l^1 y_0^t + G_2^1 y_l^{t-1}) \\
a_l^t &= g_l^t \odot \left(V_l y_0^t + U_l y_l^{t-1} + A_l^t y_l^{t-1} \right) \\
z_l^t &= \phi_z^l(a_l^t)
\end{align}

\subsubsection{Output Prediction}
\begin{itemize}
    \item $z_l^{t-1}$: the previous layer 
    \item $\phi_{o}^l$ : activation function for the prediction layer
    \item $W_l$: weights to reshape the layer before activation. 
    \item $z_{l-1,o}^t$: \textbf{prediction of the next layer. }
\end{itemize}
\begin{align}
z_{l-1,o}^t = \phi_{o}^l(W_lz_l^{t-1})
\end{align}

\subsubsection{Prediction Error}
Based on the Forward Pass, the prediction error was calculated by finding the difference between the predicted layers and their expected values. The output error was also calculated by finding the difference between the output of the first layer and the true target.
\begin{itemize}
    \item $z_{l,o}^t$: the actually outputted layer
    \item $z_l^t$ : the expected layer (not on answer, but just make sure it match from previous state)
    \item $\bar{x}$: the original input (target value)
    \item $e_1$: The error for prediction layer, to the expected value. 
    \item $e_0$ The error for the final output to the original input (true value)
\end{itemize}
 %z_{0,o}^t &= \phi_{o}^1(W_1z_1^{t-1})
\begin{align}
 z_{1,o}^t &= \phi_{o}^2(W_2z_2^{t-1}) \\
 e_1 &= z_{1,o}^t - z_1^t      %include a masking funciton
 \end {align}
Personal Notes: Could include the Masking mainly just include the state if its 1 or 0 
\begin{align}
z_{0,o}^t &= \phi_{o}^1(W_1z_1^{t-1}) \\
 e_o &= z_{0,o}^t - \bar{x}    %include a masking funciton
\end {align}

\subsection{Error Feedback}
The purpose of local learning across layers is to ensure that each layer's output aligns with both the preceding and succeeding layers. Additionally means applying an error correction for an accurate representation of each layers. 

\subsubsection{Fast weights implementation}
\hspace{.5cm}Fast weights store the connections between past and current activations—specifically, the previous and current state matrices of Layers 1 and 2. The outer product of these states produces a matrix that captures how strongly each pair of neurons co-activated. The matrix is used to update the weights for Layers 1 and 2. 
\begin{itemize}
    \item $\lambda$: decay factor, how much of the previous activity should it remember. 
    \item $\eta$: learning constant, how much of the 
\end{itemize}
This equation will incorporate the \textbf{outer product of the previous state with the current state.} \\
\begin{itemize}
    \item $A^t_2$: \textbf{Fast weight matrix — stores short-term associations between previous and current states}
    \item $z^{t-1}_2$: Pre-synaptic activity: t-1 "Caused the activation"
    \item $(z^t_2)^T$: Post-synaptic activity: t "Was the Activation"
\end{itemize}
Strengthen the fast-weight matrix [connection between neurons]. If the presynaptic and postsynaptic weights are high, then the outer product of the connection matrix will be high. \\
If neuron t is active, then neuron t+1 likely to activate again. \textbf{Stronger update, bigger influence in the next step. }
\begin{align}
\Delta A_2^t &= \eta z^{t-1}_2(z^t_2)^T \\
A_2^t &= \lambda A^{t-1}_2 + \eta z^{t-1}_2(z^t_2)^T \\
\end{align}

\subsubsection{State Correction} % this is based on the "Calculation Formula": 3.4 Error feedback and 3.5 Corrected States
Following fastweight readjustment, the state layer is adjusted based on the accumulated error from the prediction error phase. The new layer, representing the corrected state, is obtained by subtracting the error from the expected layer.
\begin{itemize}
    \item $d_l$: the error feedback produced by multiplying the error variable by the weighted Matrix
    \item $e_{l-1}$ the error calculated from the "Error Feedback stage"
    \item $E_l$: The weights for the error calculations.
    \item $\alpha$: the error calling factor (default: 0.0001)
\end{itemize}
\begin{align}
d_2 &= E_2 \cdot e_1 \\
d_1 &= E_1 \cdot e_0\\ %- e_1 \cdot \alpha 
d_l &= E_l \cdot e_{l-1}
\end{align}
\begin{itemize}
    \item $\beta$: the error feedback scaling factor (default: 0.0001)
    \item $z^t_l$ : the expected layer (not on answer, but just make sure it match from previous state) % This is representation of output: the expected layer based from the "output prediction"
    \item $\phi^l_z$ activation function for error feedback
    \item $y_l^{t}$: the new state of layer, learning form the feedback
\end{itemize}
\begin{align}
y_l^{t} &= \phi^l_z(z^t_l - \beta d_l)
\end{align}
Once we obtain the updated state and fast weights, we apply it into the forward pass of the next time step to get a more accurate output prediction:  \begin{math}( A_l^t, y_l^{t}) \end{math}
\begin{align}
a_l^{t+1} &= g_l^{t+1} \odot \left(V_l y_0^{t+1} + U_l y_l^{t} + A_l^t y_l^{t} \right)
\end{align}

\subsection{Learning Phase}
\hspace*{.5cm}Instead of backpropagation where we calculate the global weights as they are associated with multiple hidden layers, we create gates, slow, and error weights for each of the 2 layer and calculating the change in weights by the Hebbian with error. 

\subsubsection{State Errors}
Hebbian with error is calculated with the outerproduct of the error and the expected output.This works because this outer product approximates local gradient of the loss with respect to the weights 
\begin{align*}
&\text{The equations below describe the derivation behind Hebbian learning with error:} \\[10pt]
&z_l = W z_{l-1}, \quad 
L = \frac{1}{2} \| z_l - \hat{z}_l \|^2, \quad 
e = W z_{l-1} - \hat{z}_l \\[10pt]
&\frac{dL}{dW} = \frac{dL}{de} \cdot \frac{de}{dW}, \quad 
L = \frac{1}{2} e^2, \quad 
\frac{dL}{de} = e \\[10pt]
&\frac{de}{dW} = z_{l-1}, \quad 
\frac{dL}{dW} = e \cdot z_{l-1}^\top \\[10pt]
&\Delta W = -\gamma \cdot \frac{dL}{dW} = -\gamma \cdot e \cdot z_{l-1}^\top
\end{align*}


The first part is to calculate the error as the difference between the the expected layer and the output layer we received from "State Correction" This equatio  applies to both Layer 1 and Layer 2. 
\begin{itemize}
    \item $e_l$: Error calculated with the difference between the corrected and the expected layers. 
    \item $y_l$: Corrected layer state after the "State Correction" phase. 
    \item $z_l$: the expected layer (not on answer, but just make sure it match from previous state)
\end{itemize}

\subsubsection{Updating the Gates Slow and Error Weights}
We set up each of the output product by multiply the error (calculated before) and multiply it expected output which is lower layers. For the second layer: \begin{math} W_2, M_2\end{math}, the error is combined with the activation of 1st layer given that the layer2 predicts layer 1. For layer 1: \begin{math} W_1, M_1\end{math} we are multiplying the variables with the input layer given that layer 1 predicts the inputs. 

\begin{itemize}
    \item $\gamma$ is the learning rate
    \item $\Delta V_l$, : weights that incorporate the bottom up signal. The purpose is so the output matches with the lower layers. 
    \item $\Delta M_l$, : weights that incorporate the Top Down signal. The purpose is so the output matches with the Upper layers. 
    \item $\Delta U_l$, :  weights that incorporate the previous time step of same layer 
    \item $(z_{l})^T$: the state of the layer
    \item $(x_t)^T $ the input layers. 
\end{itemize}

For slow weights:
\begin{align}
\Delta V_1 &= -\gamma \cdot e_1 \cdot (x_t)^T \quad &
\Delta V_2 &= -\gamma \cdot e_2 \cdot (z_{1})^T \\
\Delta U_1 &= -\gamma \cdot e_1 \cdot (z_1^{t-1})^T \quad &
\Delta U_2 &= -\gamma \cdot e_2 \cdot (z_2^{t-1})^T\\    % z_{f,2,t-1}^T 
\Delta M_1 &= -\gamma \cdot e_1 \cdot (x_t)^T \quad &
\Delta M_2 &= -\gamma \cdot e_2 \cdot (z_{1})^T
\end{align}

 The same equations are applied for the gated weights and the error weights as well. 
For gate weights:   % I have no idea how to update the Gates in terms of calculations
\begin{align}
\Delta G_{1}^l &= -\gamma \cdot e_1 \cdot z_{f,1,t-1}^Tdo \quad &
\Delta G_{2}^l &= -\gamma \cdot e_1 \cdot x_t^T \\
\end{align}

For error feedback weights:
\begin{align}
\Delta E_1 &= -\gamma \cdot e_1 \cdot x_t^T \\
\Delta E_2 &= -\gamma \cdot e_2 \cdot z_{f,1}^T
\end{align}












%------endpoint for Model Representation
\section{Results}

\subsection{Hyperparameter Configuration}
To optimize training efficiency while maintaining model performance, we made several strategic adjustments to the hyperparameters:

\begin{itemize}
    \item \textbf{Batch Size Adjustments}:
    \begin{itemize}
        \item Training batch size (mb): Increased from 50 to 200
        \item Validation batch size (v\_mb): Increased from 100 to 400
    \end{itemize}
    These changes significantly reduced the number of iterations per epoch while maintaining stable training.

    \item \textbf{Model Architecture}:
    \begin{itemize}
        \item Hidden dimension (hid\_dim): Reduced from 1000 to 250
        The reduction in model capacity helped decrease computational requirements while still capturing essential patterns in the data.
    \end{itemize}

    \item \textbf{Optimization Parameters}:
    \begin{itemize}
        \item Learning rate (alpha): Increased from 0.075 to 0.1
        \item Momentum: Decreased from 0.95 to 0.9
        These adjustments allowed for faster convergence while maintaining training stability.
    \end{itemize}

    \item \textbf{Training Configuration}:
    \begin{itemize}
        \item Number of epochs (n\_e): Set to 30
        \item Evaluation interval (eval\_iter): Set to infinity to disable intermediate evaluations
        The configuration focused on end-to-end training without the overhead of frequent evaluations.
    \end{itemize}
\end{itemize}

These hyperparameter modifications resulted in a significant reduction in training time while maintaining competitive model performance. The changes were carefully chosen to balance computational efficiency with model effectiveness.

\subsection{Training Performance}
The model was trained for 30 epochs on the Penn TreeBank character-level dataset. The training process showed significant improvement in both training and validation metrics. Here are the key observations:

\begin{itemize}
    \item \textbf{Initial Performance}: The model started with a training loss of 3.89 and validation loss of 3.89, with very low accuracy (1.9\% training, 1.8\% validation). This low initial accuracy is expected for character-level prediction tasks, as the model must learn to predict the next character from a vocabulary of 50 characters.
    
    \item \textbf{Rapid Initial Improvement}: The model showed dramatic improvement in the first iteration, with accuracy jumping from 1.9\% to 16.6\%. It indicates that the model quickly learned basic patterns in the character sequences.
    
    \item \textbf{Convergence Pattern}: The model exhibited a non-monotonic convergence pattern:
    \begin{itemize}
        \item Early epochs (0-5): Rapid improvement followed by some fluctuation
        \item Middle epochs (6-15): Gradual improvement with occasional dips
        \item Later epochs (16-30): More stable improvement with final convergence
    \end{itemize}
    
    \item \textbf{Final Performance}: After 30 epochs, the model achieved:
    \begin{itemize}
        \item Training Loss: 2.64 (reduced from initial 3.89)
        \item Validation Loss: 2.64 (reduced from initial 3.89)
        \item Training Accuracy: 28.4\% (improved from initial 1.9\%)
        \item Validation Accuracy: 28.6\% (improved from initial 1.8\%)
        \item Perplexity: 3.81 (reduced from initial 5.61)
    \end{itemize}
    
    The accuracy is calculated at the character level, where the model predicts the next character in the sequence. For each position in the sequence, the model outputs a probability distribution over all possible characters (vocabulary size of 50). The accuracy is computed as the ratio of correctly predicted characters to the total number of valid tokens in the sequence. The final accuracy of 28.4\% indicates that the model correctly predicts the next character more than one-quarter of the time, which is significant given the complexity of natural language patterns and the large vocabulary size.
\end{itemize}

\begin{table}[h]
\centering
\caption{Training and Validation Metrics at Key Epochs}
\begin{tabular}{lcccccc}
\toprule
Epoch & Train Loss & Val Loss & Train Acc & Val Acc & Train PPL & Val PPL \\
\midrule
Initial & 3.89 & 3.89 & 1.9\% & 1.8\% & 5.61 & 5.61 \\
0 & 3.41 & 3.41 & 16.6\% & 16.6\% & 4.91 & 4.92 \\
5 & 2.88 & 2.88 & 17.8\% & 17.9\% & 4.16 & 4.16 \\
10 & 2.83 & 2.84 & 21.5\% & 21.6\% & 4.09 & 4.10 \\
15 & 2.73 & 2.73 & 23.8\% & 24.0\% & 3.94 & 3.94 \\
20 & 2.78 & 2.78 & 22.5\% & 22.7\% & 4.01 & 4.01 \\
25 & 2.69 & 2.70 & 27.8\% & 27.9\% & 3.88 & 3.89 \\
30 & 2.64 & 2.64 & 28.4\% & 28.6\% & 3.81 & 3.81 \\
\bottomrule
\end{tabular}
\label{tab:training_metrics}
\end{table}

The training progress is visualized in Figures \ref{fig:accuracy_curves_FW} through \ref{fig:parameter_norms_FW}. Figure \ref{fig:accuracy_curves_FW} illustrates the improvement in prediction accuracy, which steadily increased from an initial 1.9\% to 28.4\% by the end of training. The perplexity curves in Figure \ref{fig:perplexity_curves_FW} show a consistent decrease, indicating better model performance in predicting the next character in the sequence. Finally, Figure \ref{fig:parameter_norms_FW} tracks the growth of key model parameters throughout training, showing controlled parameter growth without any signs of instability.

\begin{figure}[h]
\centering
% \begin{minipage}{0.48\textwidth}
% \centering
% \includegraphics[width=\textwidth]{figures/loss_curves_FW.png}
% \caption{Training and Validation Loss Curves showing the convergence of the model over 30 epochs. The close alignment between training and validation curves indicates minimal overfitting.}
% \label{fig:loss_curves_FW}
% \end{minipage}
% \hfill
\begin{minipage}{0.48\textwidth}
\centering
\includegraphics[width=\textwidth]{figures/accuracy_curves_FW.png}
\caption{Training and Validation Accuracy Curves demonstrating the model's improvement in prediction accuracy from 1.9\% to 28.4\% over the training period.}
\label{fig:accuracy_curves_FW}
\end{minipage}
\end{figure}

\begin{figure}[h]
\centering
\begin{minipage}{0.48\textwidth}
\centering
\includegraphics[width=\textwidth]{figures/perplexity_curves_FW.png}
\caption{Training and Validation Perplexity Curves showing the model's improving ability to predict the next character in the sequence. Lower perplexity indicates better performance.}
\label{fig:perplexity_curves_FW}
\end{minipage}
\hfill
\begin{minipage}{0.48\textwidth}
\centering
\includegraphics[width=\textwidth]{figures/parameter_norms_FW.png}
\caption{Parameter Norm Growth Over Training for key model parameters (W1, W2, V1, V2). The controlled growth indicates stable learning without parameter explosion.}
\label{fig:parameter_norms_FW}
\end{minipage}
\end{figure}

\subsection{Model Stability}
The training process demonstrated good stability, as evidenced by:

\begin{itemize}
    \item \textbf{Consistent Improvement}: Despite some fluctuations, the overall trend showed steady improvement in both training and validation metrics.
    
    \item \textbf{Minimal Overfitting}: The gap between training and validation metrics remained small throughout training, indicating good generalization. For example, at epoch 30, the training accuracy (28.4\%) and validation accuracy (28.6\%) were very close.
    
    \item \textbf{Training Time Stability}: The training time per epoch was relatively stable, ranging from 2.5 to 3.5 hours, with later epochs showing slightly faster training times.
\end{itemize}

\subsection{Parameter Analysis}
Analysis of the parameter norms reveals interesting patterns:

\begin{itemize}
    \item \textbf{Fast Weights}: The fast weight matrices (A1, A2) maintained very small norms (around 0.00045), indicating they were effectively constrained and didn't dominate the learning process. This suggests the fast weight mechanism was working as intended, providing short-term memory without overwhelming the main network.
    
    \item \textbf{Gate Parameters}: The gate parameters showed significant growth:
    \begin{itemize}
        \item G1\_1: Grew from 2.76 to 2853.92
        \item G1\_2: Grew from 18.82 to 4953.75
        \item G2\_1: Grew from 2.77 to 3548.13
        \item G2\_2: Grew from 18.76 to 4953.75
    \end{itemize}
    This substantial growth in gate parameters indicates they played a crucial role in controlling information flow through the network.
    
    \item \textbf{Core Weights}: The main weight matrices showed controlled growth:
    \begin{itemize}
        \item W1: Grew from 5.50 to 15.34
        \item W2: Grew from 12.53 to 542.80
        \item V1: Grew from 12.46 to 210.27
        \item V2: Grew from 12.46 to 4953.86
    \end{itemize}
    The growth in these parameters was more moderate compared to the gate parameters, suggesting stable learning dynamics.
\end{itemize}

\subsection{Training Efficiency}
The training process was computationally efficient:

\begin{itemize}
    
    \item \textbf{Memory Usage}: The model maintained stable memory usage throughout training, with no signs of memory leaks or excessive resource consumption.
    
    \item \textbf{Convergence Rate}: The model showed good convergence properties:
    \begin{itemize}
        \item Initial rapid improvement in the first iteration
        \item Steady progress through middle epochs
        \item Final convergence in later epochs
    \end{itemize}
\end{itemize}


\subsection{- Comparison with LSTM Model}
The comparison between our PTNCN with fast weight and gate mechanism and the LSTM model reveals several key differences in performance and behavior:

\begin{table}[h]
\centering
\caption{Performance Comparison between PTNCN-FW and LSTM Models}
\begin{tabular}{lcccccc}
\toprule
Model & Train Loss & Val Loss & Train Acc & Val Acc & Train PPL & Val PPL \\
\midrule
PTNCN-FW (Initial) & 3.89 & 3.89 & 1.9\% & 1.85\% & 5.61 & 5.61 \\
PTNCN-FW (Final) & 2.64 & 2.64 & 28.4\% & 28.6\% & 3.81 & 3.81 \\
LSTM (Initial) & 9.94 & 9.97 & 0.72\% & 0.68\% & 14.34 & 14.38 \\
LSTM (Final) & 0.007 & 0.074 & 22.5\% & 22.7\% & 0.01 & 0.11 \\
\bottomrule
\end{tabular}
\label{tab:model_comparison}
\end{table}

\begin{figure}[h]
\centering
\begin{minipage}{0.48\textwidth}
\centering
\includegraphics[width=\textwidth]{figures/loss_curves_FW.png}
\caption{PTNCN-FW Loss Curves}
\label{fig:ptncn_loss}
\end{minipage}
\hfill
\begin{minipage}{0.48\textwidth}
\centering
\includegraphics[width=\textwidth]{figures/loss_curves_lstm.png}
\caption{LSTM Loss Curves}
\label{fig:lstm_loss}
\end{minipage}
\end{figure}

\begin{figure}[h]
\centering
\begin{minipage}{0.48\textwidth}
\centering
\includegraphics[width=\textwidth]{figures/accuracy_curves_FW.png}
\caption{PTNCN-FW Accuracy Curves}
\label{fig:ptncn_acc}
\end{minipage}
\hfill
\begin{minipage}{0.48\textwidth}
\centering
\includegraphics[width=\textwidth]{figures/accuracy_curves_lstm.png}
\caption{LSTM Accuracy Curves}
\label{fig:lstm_acc}
\end{minipage}
\end{figure}

\begin{figure}[h]
\centering
\begin{minipage}{0.48\textwidth}
\centering
\includegraphics[width=\textwidth]{figures/perplexity_curves_FW.png}
\caption{PTNCN-FW Perplexity Curves}
\label{fig:ptncn_ppl}
\end{minipage}
\hfill
\begin{minipage}{0.48\textwidth}
\centering
\includegraphics[width=\textwidth]{figures/perplexity_curves_lstm.png}
\caption{LSTM Perplexity Curves}
\label{fig:lstm_ppl}
\end{minipage}
\end{figure}

\begin{figure}[h]
\centering
\begin{minipage}{0.48\textwidth}
\centering
\includegraphics[width=\textwidth]{figures/parameter_norms_FW.png}
\caption{PTNCN-FW Parameter Norms}
\label{fig:ptncn_norms}
\end{minipage}
\hfill
\begin{minipage}{0.48\textwidth}
\centering
\includegraphics[width=\textwidth]{figures/parameter_norms_lstm.png}
\caption{LSTM Parameter Norms}
\label{fig:lstm_norms}
\end{minipage}
\end{figure}

The comparison reveals several key differences between the PTNCN-FW and LSTM models:

\begin{itemize}
    \item \textbf{Initial Performance}:
    \begin{itemize}
        \item PTNCN-FW started with better initial accuracy (1.9\% training, 1.85\% validation)
        \item LSTM had lower initial accuracy (0.72\% training, 0.68\% validation)
        \item PTNCN-FW's better initialization suggests more effective architecture design
    \end{itemize}

    \item \textbf{Convergence Pattern}:
    \begin{itemize}
        \item PTNCN-FW showed stable, monotonic improvement in accuracy
        \item LSTM exhibited more volatile behavior with frequent oscillations
        \item PTNCN-FW's smoother convergence indicates better learning dynamics
    \end{itemize}

    \item \textbf{Final Performance}:
    \begin{itemize}
        \item PTNCN-FW achieved higher final accuracy (28.4\% training, 28.6\% validation)
        \item LSTM reached lower accuracy (22.5\% training, 22.7\% validation)
        \item The 6\% accuracy advantage of PTNCN-FW demonstrates its superior performance
    \end{itemize}

    \item \textbf{Loss and Perplexity}:
    \begin{itemize}
        \item PTNCN-FW maintained consistent positive loss values (2.64 final training loss)
        \item LSTM showed highly volatile loss values, including negative values (e.g., -1.84 at epoch 0)
        \begin{itemize}
            \item The negative loss values in LSTM occur due to its implementation of sparse categorical crossentropy loss with logits, where the loss can become negative when the model's predictions are very confident
        \end{itemize}
        \begin{itemize}
            \item The LSTM's loss calculation in the training loop uses tf.keras.losses.sparse\_categorical\_crossentropy with from\_logits=True, which can produce negative values when the model's logits are strongly aligned with the correct class.
        \end{itemize}
        \item This difference in loss behavior suggests that while the LSTM might appear to have lower loss values, the negative values indicate a different loss function implementation rather than better performance
    \end{itemize}

    \item \textbf{Training Stability}:
    \begin{itemize}
        \item PTNCN-FW showed minimal gap between training and validation metrics
        \item LSTM exhibited larger fluctuations between training and validation
        \item PTNCN-FW's stability suggests better generalization capabilities
    \end{itemize}

    \item \textbf{Parameter Growth}:
    \begin{itemize}
        \item PTNCN-FW showed controlled parameter growth with fast weights maintaining small norms
        \item LSTM demonstrated more aggressive parameter growth in certain components
        \item PTNCN-FW's controlled growth indicates better regularization properties
    \end{itemize}
\end{itemize}

The comparison demonstrates that the PTNCN architecture, enhanced with fast weights and gates, offers significant advantages over traditional LSTM models for character-level language modeling tasks. The combination of local representation alignment, fast weight memory, and gating mechanisms provides a more effective approach to capturing sequential patterns in the data, resulting in better accuracy and more stable training.


\subsection{- Comparison with PTNCN Model}
The comparison between the original PTNCN model and the enhanced PTNCN model with fast weight and gate mechanisms reveals significant differences in performance and learning dynamics:

\begin{table}[h]
\centering
\caption{Performance Comparison between PTNCN and PTNCN-FW Models}
\begin{tabular}{lcccccc}
\toprule
Model & Train Loss & Val Loss & Train Acc & Val Acc & Train PPL & Val PPL \\
\midrule
PTNCN (Initial) & 3.92 & 3.92 & 1.72\% & 1.73\% & 5.65 & 5.65 \\
PTNCN (Final) & 2.42 & 2.42 & 32.77\% & 32.85\% & 3.50 & 3.50 \\
PTNCN-FW (Initial) & 3.89 & 3.89 & 1.91\% & 1.85\% & 5.61 & 5.61 \\
PTNCN-FW (Final) & 2.64 & 2.64 & 28.38\% & 28.63\% & 3.81 & 3.81 \\
\bottomrule
\end{tabular}
\label{tab:ptncn_comparison}
\end{table}

\begin{figure}[h]
\centering
\begin{minipage}{0.48\textwidth}
\centering
\includegraphics[width=\textwidth]{figures/loss_curves_ptncn.png}
\caption{PTNCN Loss Curves}
\label{fig:ptncn_loss}
\end{minipage}
\hfill
\begin{minipage}{0.48\textwidth}
\centering
\includegraphics[width=\textwidth]{figures/loss_curves_FW.png}
\caption{PTNCN-FW Loss Curves}
\label{fig:ptncn_fw_loss}
\end{minipage}
\end{figure}

\begin{figure}[h]
\centering
\begin{minipage}{0.48\textwidth}
\centering
\includegraphics[width=\textwidth]{figures/accuracy_curves_ptncn.png}
\caption{PTNCN Accuracy Curves}
\label{fig:ptncn_acc}
\end{minipage}
\hfill
\begin{minipage}{0.48\textwidth}
\centering
\includegraphics[width=\textwidth]{figures/accuracy_curves_FW.png}
\caption{PTNCN-FW Accuracy Curves}
\label{fig:ptncn_fw_acc}
\end{minipage}
\end{figure}

The comparison reveals several key differences between the original PTNCN and the enhanced PTNCN-FW models:

\begin{itemize}
    \item \textbf{Initial Performance}:
    \begin{itemize}
        \item Both models started with similar initial accuracy (PTNCN: 1.72\%, PTNCN-FW: 1.91\%)
        \item Initial loss values were comparable (PTNCN: 3.92, PTNCN-FW: 3.89)
        \item Initial perplexity values were similar (PTNCN: 5.65, PTNCN-FW: 5.61)
    \end{itemize}

    \item \textbf{Convergence Pattern}:
    \begin{itemize}
        \item PTNCN showed more rapid initial improvement in accuracy
        \item PTNCN-FW exhibited more stable learning with fewer fluctuations
        \item PTNCN reached higher peak accuracy earlier in training
    \end{itemize}

    \item \textbf{Final Performance}:
    \begin{itemize}
        \item PTNCN achieved higher final accuracy (32.77\% training, 32.85\% validation)
        \item PTNCN-FW reached lower but still competitive accuracy (28.38\% training, 28.63\% validation)
        \item PTNCN achieved lower final loss (2.42 vs 2.64) and perplexity (3.50 vs 3.81)
    \end{itemize}

    \item \textbf{Training Stability}:
    \begin{itemize}
        \item PTNCN-FW showed more consistent improvement with fewer oscillations
        \item PTNCN exhibited more volatile behavior but achieved higher peak performance
        \item Both models maintained good generalization with small training-validation gaps
    \end{itemize}

    \item \textbf{Parameter Growth}:
    \begin{itemize}
        \item PTNCN-FW showed more controlled parameter growth due to gate mechanisms
        \item PTNCN exhibited more aggressive parameter growth in core components
        \item Fast weights in PTNCN-FW maintained very small norms (around 0.00045)
        \item Gate parameters in PTNCN-FW showed significant growth (G1\_1: 2.76 to 2853.92)
    \end{itemize}

    \item \textbf{Learning Dynamics}:
    \begin{itemize}
        \item PTNCN-FW's gate mechanism provided better control over information flow
        \item Fast weights in PTNCN-FW enabled better short-term memory capabilities
        \item PTNCN showed more direct learning of patterns without intermediate control
    \end{itemize}
\end{itemize}

The comparison demonstrates that while the original PTNCN model achieved higher final accuracy, the enhanced PTNCN-FW model offers several advantages:

\begin{itemize}
    \item More stable training process with fewer oscillations
    \item Better controlled parameter growth through gate mechanisms
    \item Enhanced short-term memory capabilities through fast weights
    \item More robust learning dynamics with explicit control over information flow
\end{itemize}

These results suggest that while the addition of fast weights and gates may slightly reduce the peak performance, it provides significant benefits in terms of training stability and model interpretability. The trade-off between performance and stability makes the PTNCN-FW model potentially more suitable for real-world applications where consistent behavior is crucial.

\section{Discussion}

The experimental results reveal an interesting trade-off between model performance and training stability in our comparison of the original PTNCN and the enhanced PTNCN-FW models. This section discusses the implications of these findings and their potential impact on practical applications.

\subsection{Performance vs. Stability Trade-off}

The original PTNCN model achieved higher final accuracy (32.77\% vs 28.38\%) but exhibited more volatile training behavior. This performance difference can be attributed to several factors:

\begin{itemize}
    \item \textbf{Direct Learning vs. Controlled Learning}:
    \begin{itemize}
        \item PTNCN's direct learning approach allows for more aggressive optimization, potentially capturing more complex patterns
        \item PTNCN-FW's gate mechanism introduces an additional layer of control, which may limit the model's ability to make rapid adjustments
        \item The fast weight mechanism in PTNCN-FW focuses on short-term memory, which might not capture long-term dependencies as effectively
    \end{itemize}

    \item \textbf{Parameter Dynamics}:
    \begin{itemize}
        \item PTNCN's parameters can grow more freely, allowing for more expressive representations
        \item PTNCN-FW's gate mechanism constrains parameter growth, potentially limiting model capacity
        \item The controlled growth in PTNCN-FW, while more stable, may prevent the model from reaching optimal parameter values
    \end{itemize}

    \item \textbf{Information Flow Control}:
    \begin{itemize}
        \item PTNCN's direct information flow enables faster learning of patterns
        \item PTNCN-FW's gated architecture provides better control but may slow down the learning process
        \item The trade-off between control and speed affects the model's ability to reach optimal performance
    \end{itemize}
\end{itemize}

\subsection{Advantages of Enhanced Architecture}

Despite the lower peak performance, the PTNCN-FW model offers several significant advantages that make it more suitable for practical applications:

\begin{itemize}
    \item \textbf{Training Stability}:
    \begin{itemize}
        \item More consistent improvement across training epochs
        \item Reduced risk of catastrophic forgetting
        \item Better handling of noisy or irregular input patterns
    \end{itemize}

    \item \textbf{Model Interpretability}:
    \begin{itemize}
        \item Gate mechanisms provide insights into information flow
        \item Fast weights offer visibility into short-term memory patterns
        \item More predictable behavior in production environments
    \end{itemize}

    \item \textbf{Resource Efficiency}:
    \begin{itemize}
        \item Controlled parameter growth reduces memory requirements
        \item More efficient use of computational resources
        \item Better scalability for larger datasets
    \end{itemize}
\end{itemize}

\subsection{Implications for Real-world Applications}

The choice between PTNCN and PTNCN-FW should be based on specific application requirements:

\begin{itemize}
    \item \textbf{When to Choose PTNCN}:
    \begin{itemize}
        \item Applications requiring maximum accuracy
        \item Scenarios with stable, well-structured input data
        \item Cases where training time is not a critical factor
    \end{itemize}

    \item \textbf{When to Choose PTNCN-FW}:
    \begin{itemize}
        \item Production environments requiring stable performance
        \item Applications with noisy or irregular input patterns
        \item Scenarios where model interpretability is important
        \item Systems with limited computational resources
    \end{itemize}
\end{itemize}

\subsection{Future Directions}

The performance-stability trade-off observed in our experiments suggests several promising directions for future research:

\begin{itemize}
    \item \textbf{Architecture Improvements}:
    \begin{itemize}
        \item Development of adaptive gate mechanisms
        \item Integration of both short-term and long-term memory capabilities
        \item Dynamic parameter growth control
    \end{itemize}

    \item \textbf{Training Strategies}:
    \begin{itemize}
        \item Curriculum learning approaches
        \item Adaptive learning rate schedules
        \item Hybrid training methods combining both architectures
    \end{itemize}

    \item \textbf{Application-specific Optimizations}:
    \begin{itemize}
        \item Domain-specific gate configurations
        \item Task-specific fast weight mechanisms
        \item Customized parameter growth strategies
    \end{itemize}
\end{itemize}

The findings from this study highlight the importance of considering both performance metrics and training stability when selecting and implementing neural network architectures. While the original PTNCN model achieves higher accuracy, the enhanced PTNCN-FW model offers a more robust and practical solution for real-world applications. Future work should focus on developing architectures that can maintain the stability benefits of PTNCN-FW while approaching the performance levels of the original PTNCN model.

\section{Conclusion}

This project has presented a comprehensive study of the Parallel Temporal Neural Coding Network (PTNCN) architecture, with a particular focus on its enhanced version incorporating fast weights and gate mechanisms. Through extensive experimentation and analysis, we have demonstrated several significant contributions to the field of neural network architectures for sequential data processing.

Our work has made several important contributions: successfully integrating fast weights and gate mechanisms into the PTNCN architecture, demonstrating the effectiveness of local representation alignment in sequential learning, and developing a novel approach to information flow control through gating mechanisms. The experimental results revealed a crucial trade-off between performance and stability: while the original PTNCN achieved higher accuracy (32.77\%), the enhanced PTNCN-FW model achieved competitive accuracy (28.38\%) with significantly better stability and training dynamics.

The findings of this study have important implications for neural network architecture design. The enhanced PTNCN architecture demonstrates that it is possible to maintain high performance while ensuring training stability and practical applicability. The fast weights effectively maintained short-term memory with minimal parameter growth, while the gate mechanisms provided better control over information flow. These improvements make the model more suitable for real-world applications where stability and interpretability are crucial.

The success of this project opens new directions for neural network architecture development and suggests potential improvements in training methodologies. The findings and methodologies developed in this work can serve as a foundation for future research in neural network architecture design and implementation, particularly in applications requiring both high performance and training stability.

\section{References}
\begin{thebibliography}{9}

\bibitem{ororbia2020continual}
A.~Ororbia, A.~Mali, C.~L. Giles, and D.~Kifer,
\textit{Continual Learning of Recurrent Neural Networks by Locally Aligning Distributed Representations},
IEEE Transactions on Neural Networks and Learning Systems, vol.~31, no.~10, pp.~4267--4278, 2020.
\url{https://doi.org/10.1109/TNNLS.2019.2953622}

\bibitem{ororbia2020continual}
Ororbia, A. G., Mali, A., Giles, C. L., \& Kifer, D. (2020).
\textit{Continual Learning of Recurrent Neural Networks by Locally Aligning Distributed Representations}.
IEEE Transactions on Neural Networks and Learning Systems, 31(10), 4267--4278.
\url{https://doi.org/10.1109/TNNLS.2019.2953622}

\bibitem{ororbia2020continual}
Ororbia, A. G., Mali, A., Giles, C. L., \& Kifer, D. (2020).
\textit{ContinualPTCNC}, GitHub Repository, 2020
Available: \url{https://github.com/AnkurMali/ContinualPTNCN/tree/main}

\bibitem{cho2014rnn}
K.~Cho, B.~van Merriënboer, Ç.~Gülçehre, D.~Bahdanau, F.~Bougares, H.~Schwenk, and Y.~Bengio,
\textit{Learning Phrase Representations using RNN Encoder--Decoder for Statistical Machine Translation},
Proceedings of the 2014 Conference on Empirical Methods in Natural Language Processing (EMNLP), 2014.
\url{https://doi.org/10.3115/v1/d14-1179}

\bibitem{ptncn2025fw}
T. ~Hair, J. ~Iwasaki, \& T. ~Le.
\textit{PTNCN\_FW}, GitHub Repository, 2025.
Available: \url{https://github.com/TKAI-LAB-Mali/PTNCN_FW}

\end{thebibliography}


\end{document}